\section{Differentiation and the contact-term correction}
\label{stconv:sec:contact}
For $k\ge1$, the ordinary function $\gamma_n^{(k)}(x)$ has a power-log
singularity of order $k+1$. Its canonical Hadamard finite part on $\stconvT$
is defined by
\begin{equation}\label{stconv:eq:Fdefinition}
 \langle F_{n,k},\phi\rangle
 =\stconvCT_{\epsilon\downarrow0}
       \int_\epsilon^1\gamma_n^{(k)}(x)\phi(x)\dd x.
\end{equation}
Here $\stconvCT$ means the coefficient independent of $\epsilon$ after the
finite expansion in negative powers of $\epsilon$ and nonconstant powers
of $\log\epsilon$ is removed. The remaining error tends to zero.
This convention uses the coordinate $x\in(0,1)$ and no finite counterterm.
It agrees with Definition~\ref{stconv:def:G} when $k=0$.

\begin{lemma}[Higher Hurwitz poles as distributions]\label{stconv:lem:higherpoles}
For each $k\ge1$, as $z\to0$,
\begin{equation}\label{stconv:eq:higherpole}
 Z(1+k-z)=\frac{(-1)^k}{k!}\frac{\delta_0^{(k)}}z
 +\sum_{j\ge0}\frac{(-1)^jz^j}{j!}
           \stconvFp_{\stconvT}[\zeta^{[j]}(k+1,x)].
\end{equation}
The finite parts in this formula have the cutoff convention
\eqref{stconv:eq:Fdefinition}.
\end{lemma}
\begin{proof}
Split $\zeta(1+k-z,x)=x^{-k-1+z}+\zeta(1+k-z,1+x)$.
Subtract the Taylor polynomial of $\phi$ at zero through degree $k$.
Its remainder gives an analytic integral near $z=0$. The subtracted
monomials give terms
$\phi^{(r)}(0)/(r!(z+r-k))$, $0\le r\le k$; only $r=k$ has a pole
at zero. Since $\langle\delta_0^{(k)},\phi\rangle=(-1)^k\phi^{(k)}(0)$,
its residue is the one stated. Expanding each regular term and each
analytic remainder is exactly the constant-term prescription for the
cutoff integrals. This proves the entire Laurent expansion.
\end{proof}

\begin{theorem}[All-index differentiation anomaly]\label{stconv:thm:anomaly}
For $n\ge0$ and $k\ge1$,
\begin{equation}\label{stconv:eq:anomaly}
 D^kG_n=F_{n,k}+(-1)^{n+1}n!e_{n+1}(k)\delta_0^{(k)}.
\end{equation}
In particular, there is no anomaly when $n\ge k$.
\end{theorem}
\begin{proof}
The meromorphic distribution family satisfies
\begin{equation}\label{stconv:eq:DZ}
 D^kZ(s)=(-1)^k(s)_kZ(s+k).
\end{equation}
This follows directly from its Fourier coefficients and the Gamma
recurrence; it also follows by repeated integration by parts in a
left half-plane and meromorphic continuation. Put $s=1-z$ and expand
\[
 (1-z)_k=k!\prod_{r=1}^k(1-z/r)
        =k!\sum_{i=0}^k(-1)^ie_i(k)z^i.
\]
On the left of~\eqref{stconv:eq:DZ}, the pole is $\delta_0^{(k)}/z$ and the
regular coefficient of $z^n$ is $D^kG_n/n!$. On the right, multiplying
the regular part of Lemma~\ref{stconv:lem:higherpoles} by the polynomial gives
$F_{n,k}/n!$: this is the finite part of the ordinary pointwise
parameter-derivative identity. Multiplying its polar part by that same
polynomial gives, in addition, the regular coefficient
$(-1)^{n+1}e_{n+1}(k)\delta_0^{(k)}$. Comparison proves the formula.
\end{proof}

The first examples are
\begin{align}
 DG_0&=F_{0,1}-\delta_0',\label{stconv:eq:anomalylow}\\
 D^2G_0&=F_{0,2}-\tfrac32\delta_0'',\notag\\
 D^2G_1&=F_{1,2}+\tfrac12\delta_0''.\notag
\end{align}
The point-supported terms disappear upon restriction to $(0,1)$, so
ordinary differentiation there cannot detect them. But convolving
$\delta_0^{(k)}$ with another distribution differentiates that other
factor everywhere. Therefore these terms cannot be discarded before
convolution, even when the desired final identity is off the singular
point.

The pointwise harmonic master identity \eqref{tower:eq:masterN} remains
valid. The new contact term governs its canonical periodic extension:
point-supported derivatives vanish on the open interval but contribute
after convolution.

\subsection{Convolutions of arbitrary differentiated Stieltjes functions}
The same argument gives a closed law for all derivative indices, not
only the zeroth Stieltjes function. Put $F_{n,0}=G_n$ and define
\[
 a_{n,k}=(-1)^{n+1}n!e_{n+1}(k)\quad(k\ge1),\qquad a_{n,0}=0.
\]
Thus $F_{n,k}=D^kG_n-a_{n,k}\delta_0^{(k)}$ when $k\ge1$.

\begin{corollary}[The full differentiated convolution calculus]
\label{stconv:cor:fulltower}
Suppose $n,m,k,l\ge0$, $r=k+l\ge1$, and the exact rule from
Theorem~\ref{stconv:thm:product} is
\[
 G_n*G_m=c_{-1}\stconvId+\sum_{j=0}^{n+m+1}c_jG_j.
\]
Then
\begin{align}
 F_{n,k}*F_{m,l}={}&\sum_{j=0}^{n+m+1}c_jD^rG_j
             -a_{n,k}D^rG_m-a_{m,l}D^rG_n\label{stconv:eq:fulltower}\\
 &+(c_{-1}+a_{n,k}a_{m,l})\delta_0^{(r)}.\notag
\end{align}
Away from the integer point this reduces to
\[
 \sum_{j=0}^{n+m+1}c_j\gamma_j^{(r)}(x)
          -a_{n,k}\gamma_m^{(r)}(x)-a_{m,l}\gamma_n^{(r)}(x).
\]
In particular, every such convolution has a finite reduction to
$\zeta^{[j]}(r+1,x)$ with $0\le j\le n+m+1$.
\end{corollary}
\begin{proof}
Multiply the two identities $F=D G-a\delta$ at their indicated
orders, interpreting the zero-order case through $a_{n,0}=0$.
The product of the first terms is $D^r(G_n*G_m)$, and
$D^r\stconvId=\delta_0^{(r)}$. Restricting off zero and applying
\eqref{tower:eq:masterN} proves the last assertion.
\end{proof}

As an explicit higher-index example, the contact term vanishes for
$F_{1,1}=DG_1$. Hence
\begin{align}
 (F_{1,1}*F_{1,1})(x)
 ={}&-2\zeta'''(3,x)-9\zeta''(3,x)\label{stconv:eq:gamma1primepair}\\
 &+[4\zeta(2)-6]\zeta'(3,x)
   +[6\zeta(2)+4\zeta(3)]\zeta(3,x).\notag
\end{align}
This follows by twice differentiating~\eqref{stconv:eq:product11} and using
\eqref{tower:eq:masterN}. It illustrates that the full tower still
closes at one spectral integer even when higher Stieltjes indices occur.

\section{A two-jet collapse for every polygamma convolution}
\label{stconv:sec:poly}
Let $H_0=0$ and $H_k=\sum_{j=1}^k1/j$ for $k\ge1$. Define
\[
 \Psi_k=\stconvFp_{\stconvT}\psi^{(k)}\qquad(k\ge0)
\]
using the same canonical cutoff convention. Since $\gamma_0=-\psi$,
Theorem~\ref{stconv:thm:anomaly} gives
\begin{equation}\label{stconv:eq:PsiD}
 \Psi_0=-G_0,\qquad
 \Psi_k=-D^kG_0-H_k\delta_0^{(k)}\quad(k\ge1).
\end{equation}

\begin{theorem}[Polygamma convolution reduction]\label{stconv:thm:polygamma}
Let $k,l\ge0$ and $r=k+l\ge1$. Then
\begin{align}
 \Psi_k*\Psi_l={}&2D^rG_1+(H_k+H_l)D^rG_0\label{stconv:eq:polyfull}\\
 &+[H_kH_l-\zeta(2)]\delta_0^{(r)}.\notag
\end{align}
Consequently, for $0<x<1$,
\begin{equation}\label{stconv:eq:polyscalar}
 (\Psi_k*\Psi_l)(x)=(-1)^{r+1}r!
 \left[2\zeta'(r+1,x)+(2H_r-H_k-H_l)\zeta(r+1,x)\right].
\end{equation}
For $k=l=0$, the appropriate separate formula is
$\Psi_0*\Psi_0=2G_1-\zeta(2)\stconvId$.
\end{theorem}
\begin{proof}
Convolving~\eqref{stconv:eq:PsiD} and using that derivatives may be placed on
either factor gives
\[
 \Psi_k*\Psi_l=D^r(G_0*G_0)
 +(H_k+H_l)D^rG_0+H_kH_l\delta_0^{(r)}.
\]
The same expression is valid when one index is zero, since its harmonic
number vanishes. Apply~\eqref{stconv:eq:product00} and
$D^r\stconvId=\delta_0^{(r)}$. Restriction to $0<x<1$ gives
$2\gamma_1^{(r)}+(H_k+H_l)\gamma_0^{(r)}$.
Now use~\eqref{tower:eq:masterN} with $n=0,1$:
\[
 \gamma_1^{(r)}=(-1)^{r+1}r![\zeta'(r+1,x)+H_r\zeta(r+1,x)],
 \quad \gamma_0^{(r)}=(-1)^rr!\zeta(r+1,x).
\]
This proves~\eqref{stconv:eq:polyscalar}.
\end{proof}

For example,
\begin{align}
 (\Psi_0*\Psi_1)(x)&=2\zeta'(2,x)+\zeta(2,x),\label{stconv:eq:polyexamples}\\
 (\Psi_1*\Psi_1)(x)&=-4\zeta'(3,x)-2\zeta(3,x),\notag\\
 (\Psi_0*\Psi_2)(x)&=-4\zeta'(3,x)-3\zeta(3,x).\notag
\end{align}
Although the last two convolutions have the same total derivative order,
they have different coefficients. Confusing $\Psi_k$ with the $k$th
derivative of $\Psi_0$ would incorrectly erase this distinction.

\subsection{A convergent trigamma identity with every subtraction visible}
For $0<x<1$, set $h=x/2$ and $f(t)=\psi'(t)$. Then
\begin{align}
 \mathcal T(x)={}&2\int_0^h\left[
       \frac{f(x-t)-f(x)+tf'(x)}{t^2}
        +\psi'(1+t)f(x-t)\right]\dd t\label{stconv:eq:trigamma-integral}\\
 &+\int_x^1f(t)f(1+x-t)\dd t
      -\frac{2f(x)}h-2f'(x)\log h\notag
\end{align}
satisfies
\begin{equation}\label{stconv:eq:trigamma-value}
 \mathcal T(x)=-4\zeta'(3,x)-2\zeta(3,x).
\end{equation}
Indeed, the shift relation $\psi'(t)=t^{-2}+\psi'(1+t)$ and symmetry
about $t=x/2$ reduce the first singular convolution integral to twice
its left half. Subtracting the constant and linear Taylor terms of
$f(x-t)$ leaves a bounded quotient. The constant terms of
$\int_\epsilon^ht^{-2}\dd t$ and
$\int_\epsilon^ht^{-1}\dd t$ are $-1/h$ and $\log h$,
respectively. This proves that~\eqref{stconv:eq:trigamma-integral} is exactly
the canonical finite-part convolution, so Theorem~\ref{stconv:thm:polygamma}
applies. Every integral in~\eqref{stconv:eq:trigamma-integral} is ordinary and
absolutely convergent.



\section{Contact terms in translated derivative products}
\label{stconv:sec:translated-contact}
For $r\ge1$ let $h_{r,m}=(-1)^m m!e_{m+1}(r)$, and put $h_{0,m}=0$.
Thus \eqref{stconv:eq:anomaly} reads
$F_{m,r}=D^rG_m+h_{r,m}\delta_0^{(r)}$.
\begin{theorem}[Translated contact law]\label{shj:contactthm}
For $r,k,m,n\ge0$, $N=r+k$, and $b=1-a$ with $0<a<1$,
\begin{align}\label{shj:generalcontact}
 &\FP\int_0^1\gamma_m^{(r)}(x)\gamma_n^{(k)}(\{x+a\})\,dx\notag\\
 &\quad=(-1)^r J_{mn}^{(N)}(a)
 +(-1)^k h_{k,n}\gamma_m^{(N)}(b)
 +(-1)^r h_{r,m}\gamma_n^{(N)}(a).
\end{align}
\end{theorem}
\begin{proof}
The local finite parts use the constant-term convention
\eqref{stconv:eq:Fdefinition}. At nonzero shift the singular supports
of the two factors are disjoint, so a partition into disjoint endpoint
neighborhoods defines the product unambiguously. Substitute the two contact
identities. Periodic integration by parts gives the first term. The
shifted delta derivative acts at $x=b$, giving
$(-1)^k h_{k,n}\gamma_m^{(N)}(b)$, and the unshifted one acts at zero,
giving $(-1)^r h_{r,m}\gamma_n^{(N)}(a)$. The product of the two
point-supported terms vanishes because their supports are distinct.
This proves the formula with the exact scale-one endpoint normalization.
\end{proof}
\subsection{Every polygamma product in first Hurwitz derivatives}
For $r,k\ge0$ define
\[
\Pi_{r,k}(a)=\FP\int_0^1\psi^{(r)}(x)\psi^{(k)}(\{x+a\})\dd x.
\]
\begin{corollary}[Polygamma closure]\label{shj:polyclosure}
For $N=r+k\ge1$,
\begin{align}
\Pi_{r,k}(a)={}&(-1)^r\big[\gamma_1^{(N)}(a)-H_r\psi^{(N)}(a)\big]\notag\\
&+(-1)^k\big[\gamma_1^{(N)}(b)-H_k\psi^{(N)}(b)\big].
\label{shj:polyclosuregamma}
\end{align}
Equivalently,
\begin{align}
\Pi_{r,k}(a)=N!\bigg\{&(-1)^{k+1}\big[\zeta^{[1]}(N+1,a)
+(H_N-H_r)\zeta(N+1,a)\big]\notag\\
&+(-1)^{r+1}\big[\zeta^{[1]}(N+1,b)
+(H_N-H_k)\zeta(N+1,b)\big]\bigg\}.
\label{shj:polyclosurezeta}
\end{align}
The case $N=0$ is \eqref{shj:psipsi}.
\end{corollary}
\begin{proof}
Take $m=n=0$ in \eqref{shj:generalcontact}, use $\gamma_0^{(r)}=-\psi^{(r)}$, $h_{r,0}=H_r$, and differentiate \eqref{shj:J00}. This proves \eqref{shj:polyclosuregamma}. The pointwise identities
\begin{align*}
\psi^{(N)}(x)&=(-1)^{N+1}N!\zeta(N+1,x),\\
\gamma_1^{(N)}(x)&=(-1)^{N+1}N!
\big[\zeta^{[1]}(N+1,x)+H_N\zeta(N+1,x)\big]
\end{align*}
then give \eqref{shj:polyclosurezeta}. The second follows directly by expanding
$\partial_x^N\zeta(1+u,x)=(-1)^N(1+u)_N\zeta(N+1+u,x)$ at $u=0$; these pointwise ladders are also part of the repository's existing development \cite{stconv:shifted-report}.
\end{proof}

Two instructive special cases are
\begin{align}
\Pi_{0,1}(a)&=\zeta^{[1]}(2,a)+\zeta(2,a)-\zeta^{[1]}(2,b)
=J_{00}'(a)+\psi'(b),\label{shj:01}\\
\Pi_{1,1}(a)&=2\zeta^{[1]}(3,a)+2\zeta^{[1]}(3,b)
+\zeta(3,a)+\zeta(3,b).\label{shj:11}
\end{align}
The term $+\psi'(b)$ in \eqref{shj:01} is an explicit counterexample to the unqualified rule ``differentiate the finite-part identity under the integral.'' It is not a counterexample to ordinary differentiation under a genuinely dominated integral.

For reproducibility, let $I_{r,k}(\epsilon)$ be the sum of the two raw cutoff integrals, with arguments split as in \eqref{shj:Jcutoff}. Then
\begin{align}
\Pi_{0,1}(a)=\lim_{\epsilon\downarrow0}\left[
I_{0,1}(\epsilon)-\frac{\psi(b)}\epsilon
-\big(\psi'(a)-\psi'(b)\big)\log\epsilon\right],\label{shj:cut01}\\
\Pi_{1,1}(a)=\lim_{\epsilon\downarrow0}\left[
I_{1,1}(\epsilon)-\frac{\psi'(a)+\psi'(b)}\epsilon
+\big(\psi''(a)+\psi''(b)\big)\log\epsilon\right].\label{shj:cut11}
\end{align}
These counterterms can be derived directly from $\psi(x)=-1/x-\gamma+O(x)$ and $\psi'(x)=x^{-2}+O(1)$. They give numerical tests of the contact law that do not implement distributional differentiation.


The full default replays are preserved under \src{verification/eighth-replay/}.
The direct cutoff polygamma calculations check the contact correction
without implementing distributional differentiation. The all-index law
rests on the Laurent residues and the proof above. The ordinary pointwise
master identity has not been retracted or altered.
